\chapter{Conclusiones y recomendación}\label{cap:concl}

\begin{enumerate}
  \item \textbf{El rotor libre cumple \Req{C4} con $R=\SI{0.20}{m}$, pero con poco margen.}
        La BEMT (cap.~\ref{cap:bemt}) da \SI{6.9}{m/s} con el drogue atado y
        \SI{7.9}{m/s} sin su aporte, frente a \SI{6.4}{m/s} de la estimación base. El
        rotor frena como un disco de $\CR\approx0{,}97$, no de 1,2. Llegar a
        \SI{5}{m/s} exactos pediría $R\approx\SI{0.27}{m}$.
  \item \textbf{El límite lo pone el área, no el perfil.} Cambiar paso, cuerda o
        torsión mueve $\Vr$ menos de \SI{0.5}{m/s}; el perfil decide sobre todo la
        velocidad de giro.
  \item \textbf{El rotor gira a unas \SI{1500}{rpm}}, no a 1150: la sección trabaja con
        $C_l\approx0{,}47$, en la rama baja de la polar (cap.~\ref{cap:polares}). La
        fuerza centrífuga sube a unos \SI{22}{N} por pala.
  \item \textbf{El paso debe ser negativo} (recomendado \ang{-3}; aceptable \ang{-2} a
        \ang{-4}): con paso $\ge\ang{0}$ el rotor no arranca solo. Con paso negativo hay
        una sola rama estable y llega al 90\,\% de su giro en unos \SI{2.3}{s}.
  \item \textbf{La estructura necesita tres cambios} (cap.~\ref{cap:estructura}): pala con
        laminado a $\pm\ang{45}$ y varilla de carbono en el borde de ataque para alejar la
        resonancia de torsión y el flameo; tope de batimiento elástico a \ang{+15}, con
        brazo $\ge\SI{10}{mm}$ (la conicidad real es de \ang{6.1}); y un anillo de espuma
        de \SI{25}{mm} en la base para el aterrizaje.
  \item \textbf{La lógica de liberación es robusta en simulación}
        (cap.~\ref{cap:electronica}): filtro de Kalman y votación 2 de 3 para el apogeo,
        con respaldo por tiempo.
  \item \textbf{La probabilidad estimada de no cumplir \Req{C4} es del 3 al 6\,\% por
        vuelo} (cap.~\ref{cap:confiabilidad}). Es una estimación de ingeniería que hay que
        actualizar con los ensayos.
  \item \textbf{El paracaídas sigue siendo la opción más segura}
        (cap.~\ref{cap:decision}); el autogiro gana cuando pesa el valor técnico del
        proyecto.
\end{enumerate}

\begin{resultado}[Recomendación]
Avanzar con la configuración A: rotor de 4 palas de placa de carbono curvada al 7\,\%
($156\times45$\,mm, laminado con varilla en el borde de ataque), paso de \ang{-3}, sin
torsión, bisagras de batimiento con tope elástico a \ang{+15} y pasador de Ø\,3, eje fijo
hueco con la cuerda del drogue por dentro, drogue de Ø\,\SI{25}{cm}, traba por anillo
interno con servo y anillo amortiguador en la base. Mantener un paracaídas de
Ø\,\SI{70}{cm} como plan~B hasta que los ensayos confirmen $\Vr$.
\end{resultado}

\section*{Decisiones abiertas}
\begin{itemize}
  \item Contenedor: confirmar con la organización si alcanza con la copa del drogue.
  \item Uso de \texttt{PROBE\_RELEASE}, \texttt{PAYLOAD\_RELEASE} y del comando
        \texttt{MEC} en vuelo.
  \item Autorización de la ANAC para soltar cargas desde un drone en los ensayos de caída
        (hacen falta al menos \SIrange{80}{100}{m}).
\end{itemize}

\section*{Próximos pasos}
\begin{enumerate}
  \item Modelo CAD completo y cierre del presupuesto de masa con los cambios de
        estructura.
  \item Fabricar palas y portapalas de \ang{-2} a \ang{-4}; ensayos E1 y E2.
  \item E3 para medir $\CR$ real y el rozamiento de los rodamientos; recalcular con la
        BEMT.
  \item Caídas desde drone (E4), vibración y choque (E5) e integración (E6).
  \item Trabajo futuro: simulación dinámica completa del descenso y análisis de
        estabilidad del conjunto rotor–cuerpo–drogue (anillo de vórtices, péndulo,
        giro del cuerpo y aletas antigiro).
\end{enumerate}
